\chapter{دوال دلتا، والحقول الكمومية، وقوانين الانحفاظ}

الحقل الكمومي البسيط الذي قُدِّم سابقًا يكفي لوصف التشتت والخلق والفناء على نحو إجمالي، لكن يمكن استخلاص المزيد منه. وعلى وجه الخصوص، يبين كيف تدخل قوانين الانحفاظ في عمليات الجسيمات. إن مقررًا كاملًا في نظرية الحقل الكمومي يستغرق وقتًا أطول بكثير من بضع محاضرات؛ أما هنا فالهدف أضيق. وسوف نستخدم أبسط حقل لنرى، على نحو آلي، من أين يأتي انحفاظ الطاقة، ولماذا لا يلزم، في وجود هدف ثابت، أن ينحفظ زخم الجسيم المتشتت، وكيف يشفّر ثبات الطور انحفاظ الشحنة.

قبل الرجوع إلى الحقل، نحتاج إلى فقرتين رياضيتين: دالة دلتا لديراك، وفعل مؤثري الخلق والفناء على متجهات البرا.

\section{دالة دلتا لديراك}

سمِّ الإحداثي \(y\)، من غير أن نقرر ما إذا كان سيمثل لاحقًا الموضع أو عدد الموجة أو شيئًا آخر. ابدأ بدالة عادية على هيئة حدبة متمركزة قرب \(y=a\). اجعل الحدبة أضيق فأضيق. ولو أُبقي ارتفاعها ثابتًا لانكمشت مساحتها إلى الصفر، لذا ارفع الارتفاع كلما نقص العرض، مع إبقاء حاصل ضرب الارتفاع في العرض مساويًا للواحد. والشيء الحدّي الناتج هو دالة دلتا لديراك \(\delta(y-a)\).

فهي تنعدم بعيدًا عن \(y=a\)، وتتركز بحدة لا نهائية عند \(y=a\)، ومساحتها تساوي واحدًا:
\begin{equation}
\delta(y-a)=0\quad (y\neq a),
\qquad
\int_{-\infty}^{\infty}dy\,\delta(y-a)=1.
\end{equation}
والنقطة المهمة هي اجتماع التوضع مع ثبات المساحة. فتضييق القمة لا يزيلها، لأن ارتفاعها يزداد في الوقت نفسه.

\begin{figure}[h]
\centering
\includegraphics[width=0.68\textwidth]{lecture_04_figure_02.png}
\caption{شوكة دالة دلتا ذات مساحة وحدوية متمركزة عند \(y=a\).}
\lectureframe{00:05:50}
\end{figure}

ويمكن تمثيل الصورة نفسها على نحو مثالي كشَوكة عمودية:
\begin{figure}[h]
\centering
\begin{tikzpicture}[scale=1.0]
\draw[->] (-3.1,0) -- (3.3,0) node[right] {$y$};
\draw[very thick] (0,0) -- (0,2.9);
\draw (0,-0.15) node[below] {$y=a$};
\node at (1.35,1.65) {$\delta(y-a)$};
\end{tikzpicture}
\caption{التموضع المثالي لـ \(\delta(y-a)\) عند صفر وسيطها.}
\end{figure}

تظهر دالة دلتا طبيعيًا في تكاملات فورييه. ولنعد لحظةً إلى فترة مكانية دورية طولها الكلي \(L\). أحد الأنماط الدورية هو
\begin{equation}
e^{ikx},
\end{equation}
حيث إن \(k\) هو أحد أعداد الموجة المسموح بها
\begin{equation}
k_n=\frac{2\pi n}{L},
\qquad n\in\mathbb Z.
\end{equation}
ويمكن تمثيل الفترة على أنها \(0\leq x<L\)، لكن من الأكثر تناظرًا وضع الأصل في المنتصف واستخدام
\begin{equation}
-\frac L2\leq x<\frac L2.
\end{equation}
وتُعَدّ النهايتان متطابقتين: فعبور \(L/2\) يعيد المرء إلى \(-L/2\).

كامِلْ أحد الأنماط المسموح بها على امتداد الفترة الدورية كلها:
\begin{equation}
I_L(k)=\int_{-L/2}^{L/2}dx\,e^{ikx}.
\end{equation}
إذا كان \(k=0\)، فالدالة تحت التكامل تساوي واحدًا على نحو مطابق والنتيجة هي \(L\). وإذا كانت \(k\neq0\) قيمةً مسموحًا بها، فإن الجزأين الموجب والسالب من التذبذب يتلاشيان على العدد الكامل من الدورات. ومن ثم
\begin{equation}
I_L(k)=
\begin{cases}
L,&k=0,\\[4pt]
0,&k\neq0,
\end{cases}
\qquad
k=\frac{2\pi n}{L}.
\end{equation}

والآن اعتبر \(I_L(k)\) دالةً على المجموعة المنفصلة من قيم \(k\) المسموح بها. والتباعد بين القيم المتجاورة هو
\begin{equation}
\Delta k=\frac{2\pi}{L}.
\end{equation}
وعند \(k=0\) يكون ارتفاع الدالة \(L\)؛ وعند كل قيمة مسموح بها أخرى تكون صفرًا. ولغرض رسم الشوكة المنفصلة، أعطها عرضًا يساوي تباعدًا واحدًا في \(k\). وعندئذ تكون مساحتها
\begin{equation}
(\text{height})(\text{width})
=L\frac{2\pi}{L}=2\pi.
\end{equation}
ومع ازدياد \(L\) يزداد الارتفاع، وينكمش التباعد، وتبقى المساحة \(2\pi\). وفي حد الحجم اللانهائي تصبح الشوكة \(2\pi\delta(k)\)، فنحصل على هوية فورييه الرسمية
\begin{equation}
\int_{-\infty}^{\infty}dx\,e^{ikx}=2\pi\delta(k).
\end{equation}
ويظهر المتغيران \(x\) و\(k\) على نحو متناظر في الأسّي. ويعطي تبديل دوريهما الهوية المرافقة
\begin{equation}
\int_{-\infty}^{\infty}dk\,e^{ikx}=2\pi\delta(x).
\end{equation}
وسرعان ما ستحوّل هاتان القاعدتان الأوليتان تكاملًا على جميع أزمنة التفاعل إلى شرط لانحفاظ الطاقة.

\section{البراوات والكيتات والمؤثرات المؤثرة إلى اليسار}

الكيت ترميز رمزي لحالة كمومية. ولغرضنا الحالي، ستُكتب الحالة الابتدائية على صورة كيت والحالة النهائية على صورة برا. ووضع برا مع كيت يعطي جداءً داخليًا، وهو عدد. وهذه الأعداد، أو مربعات مقاديرها المطلقة، هي في النهاية ما يمكن مقارنته بالتجربة.

\begin{classroomqa}
\lecturetimestamp{00:15:30}
\audiencequestion{هل كان قد جرى بالفعل تقديم كلمة ``ket'' من قبل؟}
\lecturerresponse{نعم. الكيت هو أحد نصفي ترميز البرا-كيت؛ أما النصف الآخر، الذي يُقدَّم الآن، فهو البرا.}
\end{classroomqa}

في المذبذب التوافقي، يمكن كتابة حالات عدد الإشغال إما على صورة كيتات \(\lvert n\rangle\) أو على صورة براوات \(\langle n\rvert\). وجداؤها الداخلي هو
\begin{equation}
\langle m\vert n\rangle=\delta_{mn}
=\begin{cases}
1,&m=n,\\
0,&m\neq n.
\end{cases}
\end{equation}
ودلتا كرونيكر \(\delta_{mn}\) هي النظير المتقطع لدالة دلتا لديراك.

تؤثر مؤثرا الخلق والفناء في الكيتات بالصورة المألوفة:
\begin{align}
a^\dagger\lvert n\rangle&=\sqrt{n+1}\,\lvert n+1\rangle,\\
a\lvert n\rangle&=\sqrt n\,\lvert n-1\rangle.
\end{align}
ويمكن أيضًا الترميز لمؤثري الخلق والفناء بـ \(a^+\) و\(a^-\)؛ وهنا سيُستعمل \(a^\dagger\) و\(a\) على امتداد ما يلي.

كيف ينبغي للمؤثرات نفسها أن تؤثر في البراوات؟ يضع الترميز المؤثر إلى يمين البرا عندما يكون فعله إلى اليسار. ويُختار التعريف بحيث تكون لعنصر المصفوفة القيمة نفسها سواء طُبِّق المؤثر أولًا على الكيت أم طُبِّق أولًا على البرا. انظر إلى
\begin{equation}
\langle n\rvert a^\dagger\lvert m\rangle.
\end{equation}
والتأثير في الكيت يعطي
\begin{align}
\langle n\rvert a^\dagger\lvert m\rangle
&=\sqrt{m+1}\,\langle n\vert m+1\rangle\\
&=\sqrt{m+1}\,\delta_{n,m+1}.
\end{align}
وهذا غير صفري فقط إذا \(n=m+1\)، أو على نحو مكافئ إذا \(m=n-1\). ولإعادة إنتاج الشرط نفسه بالتأثير في البرا، يجب على \(a^\dagger\) أن يُنقص مؤشرها:
\begin{equation}
\langle n\rvert a^\dagger
=\sqrt n\,\langle n-1\rvert.
\end{equation}
والعامل العددي هو \(\sqrt n\)، لأن \(n=m+1\) على عنصر المصفوفة غير الصفري. ويلتزم مؤثر الفناء بالقاعدة المكملة
\begin{equation}
\langle n\rvert a
=\sqrt{n+1}\,\langle n+1\rvert.
\end{equation}
وهكذا فإن \(a^\dagger\) يرفع كيتًا لكنه يُنقص برا، في حين أن \(a\) يُنقص كيتًا لكنه يرفع برا.

وتظهر الفائدة في عنصر مصفوفة مؤثر العدد
\begin{equation}
\langle n\rvert a^\dagger a\lvert n\rangle.
\end{equation}
فبالتأثير إلى اليمين،
\begin{align}
a\lvert n\rangle&=\sqrt n\,\lvert n-1\rangle,\\
a^\dagger\lvert n-1\rangle&=\sqrt n\,\lvert n\rangle,
\end{align}
ومن ثم
\begin{equation}
\langle n\rvert a^\dagger a\lvert n\rangle
=n\langle n\vert n\rangle=n.
\end{equation}
أما بالتأثير إلى اليسار بدلًا من ذلك،
\begin{align}
\langle n\rvert a^\dagger&=\sqrt n\,\langle n-1\rvert,\\
\langle n-1\rvert a&=\sqrt n\,\langle n\rvert,
\end{align}
فنحصل على الجواب نفسه. وليس من الضروري تقرير أي الاتجاهين أكثر أساسية؛ فقد نُظِّم الحساب بحيث يعطي أيٌّ من الاتجاهين العدد نفسه.

\section{حقل كمومي في المكان والزمان}

والسبب في هذا التنظيم الحسابي هو فيزياء الجسيمات. فعلى المقاييس المجهرية، تكون الوسيلة التجريبية الرئيسية هي تصادم الجسيمات ورصد ما يخرج منه. وقد تتحول حالة ابتدائية فيها جسيم واحد أو جسيمان إلى حالة فيها أربعة أو خمسة أو سبعة أو تسعة أو أي عدد آخر. وتوفّر الحقول الكمومية، المبنية من مؤثري الخلق والفناء، لغةً لمثل هذه العمليات.

فلنبدأ ببعد مكاني واحد. ولا شيء جوهريًا يمنع استعمال ثلاثة أبعاد: إذ يصبح كل من \(x\) و\(k\) ببساطة متجهًا ذا ثلاث مركبات. ومن الملائم أيضًا استعمال وحدات يكون فيها
\begin{equation}
\hbar=1.
\end{equation}
وعندئذ تكون للزخم ومتجه الموجة القيمة العددية نفسها،
\begin{equation}
\vec p=\vec k.
\end{equation}

\begin{figure}[h]
\centering
\includegraphics[width=0.58\textwidth]{lecture_04_figure_03.png}
\caption{الوحدات الطبيعية \(\hbar=1\) والتطابق \(\vec p=\vec k\).}
\lectureframe{00:36:56}
\end{figure}

اعتمد الاصطلاح القائل إن \(\psi\) يفني و\(\psi^\dagger\) يخلق:
\begin{align}
\psi(x)&=\sum_k a(k)e^{ikx},\\
\psi^\dagger(x)&=\sum_k a^\dagger(k)e^{-ikx}.
\end{align}
ويكون الجمع على \(k=2\pi n/L\) في الفضاء الدوري، ويصير مستمرًا في الفضاء اللانهائي. والمعادلة الثانية هي المرافق الهيرميتي للأولى. ويمكن أيضًا النظر إليها بوصفها النظير المؤثري للمرافق المركب.\footnote{تُعنوِن السبورة مؤقتًا توسع حقل الخلق \(\sum_k a^\dagger(k)e^{-ikx}\) بالرمز \(\Psi\)، ثم تتردد لاحقًا بشأن إسناد الخنجر. وهنا ثُبِّتت التسميات على نحو عام بحيث إن \(\psi\) يفني و\(\psi^\dagger\) يخلق، تمامًا كما في حساب التشتت اللاحق.}

لماذا هذا التعريف؟ أفضل جواب في هذه المرحلة هو أنه مفيد. فهو مجموع بسيط على أنماط الزخم، ومعاملاته هي على وجه الدقة المؤثرات التي تزيل الجسيمات أو تضيفها. أما التبرير التفصيلي فيأتي من الديناميكيات ومن الحسابات التي يتيحها.

\editorialnote{مقطع قصير من المصدر يأتي مباشرةً بعد تقديم الحقل المكاني المرافق غير واضح بما يكفي ليستند إليه بيان إضافي عن التطبيع أو التفسير. ويستأنف المحتوى الموثوق عند فائدة التعريف.}

حتى الآن تعتمد الحقول على الموضع لا على الزمن. فالحقول تتذبذب، ولكل عدد موجي \(k\) تردد مرتبط به هو \(\omega(k)\). ولشيء ذي تردد محدد \(\omega\) تبعية زمنية
\begin{equation}
e^{i\omega t}=\cos(\omega t)+i\sin(\omega t).
\end{equation}
ألحق الطورين المترافقين بتوسعي الحقلين:
\begin{align}
\psi(x,t)&=\sum_k a(k)e^{ikx}e^{-i\omega(k)t},\\
\psi^\dagger(x,t)&=\sum_k a^\dagger(k)e^{-ikx}e^{i\omega(k)t}.
\end{align}
والآن يتذبذب كل نمط مكاني عند التردد الملائم لعدد موجته، ويعتمد الحقل على المكان والزمان معًا.

\section{معادلة الحقل غير النسبية}

العلاقة بين \(\omega\) و\(k\) هي التي تحدد المعادلة التفاضلية التي يحققها الحقل. وبالنسبة إلى كمٍّ واحد، فإن علاقات الموجة-الجسيم هي
\begin{equation}
E=\hbar\omega,
\qquad
p=\hbar k.
\end{equation}

\begin{figure}[h]
\centering
\includegraphics[width=0.62\textwidth]{lecture_04_figure_04.png}
\caption{علاقتا الطاقة-التردد والزخم-عدد الموجة \(E=\hbar\omega\) و\(p=\hbar k\).}
\lectureframe{00:41:28}
\end{figure}

عند اتخاذ \(\hbar=1\)، تصبح الطاقة هي التردد ويصبح الزخم هو عدد الموجة:
\begin{equation}
E=\omega,
\qquad
p=k.
\end{equation}
وبالنسبة إلى جسيم غير نسبي يتحرك أبطأ بكثير من الضوء، فإن الطاقة الحركية والزخم يحققان
\begin{equation}
E=\frac{p^2}{2m}.
\end{equation}
وهذه هي العلاقة نفسها المستحصلة من \(E=\tfrac12 mv^2\) و\(p=mv\). وبجمع هذه العلاقات نحصل على علاقة التشتت غير النسبية
\begin{equation}
\omega(k)=\frac{k^2}{2m}.
\end{equation}

أما حقل الخلق، فتوسع الأنماط له هو
\begin{equation}
\psi^\dagger(x,t)=\sum_k a^\dagger(k)e^{-ikx}e^{i\omega(k)t}.
\end{equation}

\begin{figure}[h]
\centering
\includegraphics[width=0.66\textwidth]{lecture_04_figure_05.png}
\caption{توسع الأنماط المعتمد لبدء اشتقاق معادلة الحقل غير النسبية.}
\lectureframe{00:43:23}
\end{figure}

اشتق بالنسبة إلى الزمن:
\begin{equation}
\dot\psi^\dagger(x,t)
=\sum_k i\omega(k)a^\dagger(k)e^{-ikx}e^{i\omega(k)t}.
\end{equation}
واشتق مرة واحدة بالنسبة إلى الموضع:
\begin{equation}
\frac{\partial\psi^\dagger}{\partial x}
=\sum_k(-ik)a^\dagger(k)e^{-ikx}e^{i\omega(k)t}.
\end{equation}
لا توجد علاقة مباشرة من الرتبة الأولى بين هاتين المشتقتين، لأن المشتقة الزمنية تعطي العامل \(\omega(k)\)، في حين أن مشتقة مكانية واحدة تعطي العامل \(k\). ولرؤية هذا الفرق صراحةً، تخيل علاقة التشتت الخطية
\begin{equation}
\omega(k)=k.
\end{equation}
عندئذٍ تعطي المشتقة الزمنية \(i\omega=ik\)، بينما تعطي المشتقة المكانية الأولى \(-ik\). ولكل نمط على حدة،
\begin{equation}
\frac{\partial\psi^\dagger}{\partial t}
=-\frac{\partial\psi^\dagger}{\partial x}.
\end{equation}
وسيحقق ذلك الحقل الافتراضي معادلة من الرتبة الأولى. أما هنا، فلدينا \(\omega=k^2/(2m)\)، ولذلك نحتاج إلى مشتقة مكانية ثانية:
\begin{align}
\frac{\partial^2\psi^\dagger}{\partial x^2}
&=\sum_k(-k^2)a^\dagger(k)e^{-ikx}e^{i\omega(k)t}\\
&=-2m\sum_k\omega(k)a^\dagger(k)e^{-ikx}e^{i\omega(k)t}.
\end{align}

\begin{classroomqa}
\lecturetimestamp{00:47:14}
\audiencequestion{هل من المفترض حقًا أن تظهر هنا إشارة السالب الناتجة عن الاشتقاقات المكانية؟}
\lecturerresponse{نعم. فكل مشتقة بالنسبة إلى \(x\) لـ \(e^{-ikx}\) تعطي العامل \(-ik\)، ومشتقتان تعطيان \((-ik)^2=-k^2\). ولذلك فإن التعويض بـ \(k^2=2m\omega\) يعطي \(-2m\omega\).}
\end{classroomqa}

إن ضرب المشتقة الزمنية في \(i\) يعطي معامل النمط نفسه الذي ينتج من قسمة المشتقة المكانية الثانية على \(2m\):
\begin{align}
i\dot\psi^\dagger
&=-\sum_k\omega(k)a^\dagger(k)e^{-ikx}e^{i\omega(k)t},\\
\frac{1}{2m}\frac{\partial^2\psi^\dagger}{\partial x^2}
&=-\sum_k\omega(k)a^\dagger(k)e^{-ikx}e^{i\omega(k)t}.
\end{align}
ومن ثم
\begin{equation}
i\frac{\partial\psi^\dagger}{\partial t}
=\frac{1}{2m}\frac{\partial^2\psi^\dagger}{\partial x^2},
\end{equation}
أو، بصورة مكافئة،
\begin{equation}
-i\frac{\partial\psi^\dagger}{\partial t}
=-\frac{1}{2m}\frac{\partial^2\psi^\dagger}{\partial x^2}.
\end{equation}

\begin{figure}[h]
\centering
\includegraphics[width=0.72\textwidth]{lecture_04_figure_06.png}
\caption{المشتقات الزمنية والمكانية لتوسعة الأنماط لحقل الإنشاء.}
\lectureframe{00:49:06}
\end{figure}

\begin{classroomqa}
\lecturetimestamp{00:49:43}
\audiencequestion{هل المعامل هو \(1/(2m)^2\)؟}
\lecturerresponse{لا. بل هو \(1/(2m)\).}
\end{classroomqa}

يحقق حقل الإفناء المرافق الصيغة المألوفة لـ Schr\"odinger
\begin{equation}
i\frac{\partial\psi}{\partial t}
=-\frac{1}{2m}\frac{\partial^2\psi}{\partial x^2}.
\end{equation}
ومعادلتا الحقل هما المرافق الهرميتي إحداهما للأخرى.\footnote{أثناء التحقق من الإشارة، ظهر على اللوح لفترة وجيزة الحد المكاني بالإشارة المعاكسة. والاشتقاق المباشر لأطوار الأنماط المعروضة يعطي زوج المعادلتين المترافقين المكتوبين هنا.}

للمعادلة الصيغة نفسها التي لمعادلة Schr\"odinger في ميكانيكا الكم الأولية، لكن الكائن الذي يحققها مختلف. ففي ميكانيكا الكم الأولية، تكون الدالة الموجية دالةً ذات قيم مركبة يُفسَّر مربع قيمتها المطلقة تفسيرًا احتماليًا. أما هنا فإن \(\psi\) و\(\psi^\dagger\) مؤثران: \(\psi\) يزيل جسيمًا و\(\psi^\dagger\) يضيف جسيمًا. ويوفر هذا الزوج المرافق صلة مباشرة بين الجسيمات والحقول.

عندما تشغل كمّات كثيرة الأنماط المعنية، يمكن للحقل أن يتصرف تقريبًا تصرفًا كلاسيكيًا، تمامًا كما يمكن لمذبذب توافقي شديد الإثارة أن يبدو كلاسيكيًا، في حين أن مذبذبًا بكم واحد لا يبدو كذلك. ويجعل الإشغال الكبير وجود مقياس كلاسيكي كبير ممكنًا، ويمكن أن يجعل التقلبات النسبية المناسبة صغيرة.

\editorialnote{إن الإشغال الكبير وحده لا يكفي للحصول على حقل كلاسيكي. فثمة حاجة أيضًا إلى حالة شبه كلاسيكية متماسكة الطور؛ أما الحالة الذاتية للعدد فلا طور محددًا لها، ولا يلزم أن يكون لها حقل متوسط غير صفري.}

\section{هدف ثابت وسعة تشتت}

كيف تُستخدم هذه الحقول؟ انظر في أبسط عملية تشتت ممكنة. يصادف جسيم هدفًا ثقيلًا إلى حد يُهمَل معه ارتداده. ويظل الهدف ثابتًا في الفضاء إلى الأبد. يمكن أن ينعكس الجسيم، أو ينفذ، أو يُعاد توجيهه؛ وفي ثلاثة أبعاد قد يختلف اتجاه خروجه عن اتجاه دخوله.

ليس من الضروري أن يكون زخم الجسيم هو نفسه قبل التصادم وبعده. فالزخم محفوظ للنظام الفيزيائي الكامل، لكن الهدف يمتص الزخم المفقود. وما إن تُحذف حركة الهدف من الوصف الرياضي حتى لا يعود زخم الجسيم وحده محفوظًا.

أما طاقته فتسلك سلوكًا مختلفًا في هذا النموذج المثالي. تخيل رمي كرة تنس على جدار. إن كرة التنس الحقيقية تسخن وتتشوه، لكن إذا أُهمِلت هذه التأثيرات فإنها تعود بالطاقة الحركية نفسها، رغم أن زخمها ينعكس. وقد صُمم نموذج الهدف الثابت لإعادة إنتاج هذا الفرق: طاقة الجسيم محفوظة، بينما قد يتغير زخمه.

ارسم الفضاء أفقيًا والزمن عموديًا. الهدف هو الخط العمودي \(x=0\). تصل إلى الهدف حالة جسيم واحد داخلة ذات زخم \(k_i\)، وتغادره حالة جسيم واحد خارجة ذات زخم \(k_f\).

\begin{figure}[h]
\centering
\begin{tikzpicture}[scale=1.0]
\draw[->] (-3.2,0) -- (3.2,0) node[right] {space};
\draw[->] (0,-0.3) -- (0,4.2) node[above] {time};
\draw[very thick] (0,0) -- (0,3.8);
\node[left] at (-0.15,3.55) {$x=0$};
\draw[thick,->] (-2.2,0.8) -- (0,2.0);
\node[below] at (-1.45,1.05) {$k_i$};
\draw[thick,->] (0,2.0) -- (2.2,3.2);
\node[above] at (1.5,2.9) {$k_f$};
\fill (0,2.0) circle (1.4pt);
\end{tikzpicture}
\caption{يتشتت جسيم ذو زخم \(k_i\) عن هدف مثبت عند \(x=0\) ويخرج بزخم \(k_f\).}
\end{figure}

استخدم لغة الامتصاص وإعادة الإصدار. فالهدف يُفني الجسيم الداخل ويُنشئ فورًا الجسيم الخارج، وربما بزخم مختلف. وهذا مجرد نموذج رياضي لحدث التشتت؛ فالذرة الحقيقية، مثلًا، يمكن أن تمتص فوتونًا، وتبقى مثارة مدةً من الزمن، ثم تُصدِره لاحقًا.

\begin{classroomqa}
\lecturetimestamp{01:02:28}
\audiencequestion{لماذا يجب امتصاص الجسيم وإعادة إصداره فورًا؟}
\lecturerresponse{لا يلزم ذلك في الطبيعة. إن الامتصاص وإعادة الإصدار الآنيين عند نقطة واحدة هما النموذج التبسيطي المستخدم هنا.}
\end{classroomqa}

تحتوي الحالة الابتدائية على جسيم واحد في النمط \(k_i\) ولا شيء في أي نمط آخر. وبلغة أعداد الإشغال تكون كل خانة صفرًا، ما عدا الخانة المرتبطة بـ \(k_i\)، فهي واحد. ويكفي اختصار القائمة الكاملة بالصيغة
\begin{equation}
\lvert k_i\rangle=a^\dagger(k_i)\lvert0\rangle.
\end{equation}
وبالمثل،
\begin{equation}
\langle k_f\rvert=\langle0\rvert a(k_f).
\end{equation}

يُمثَّل الإفناء عند الهدف وعند الزمن \(t\) بـ \(\psi(0,t)\)؛ ويُمثَّل الإنشاء الفوري من جديد عند النقطة نفسها والزمن نفسه بـ \(\psi^\dagger(0,t)\). ولا يوجد زمن مفضل يقع عنده الحدث. إذ يمكن أن يحدث كلما وصل الجسيم الداخل إلى الهدف، وبشدة التفاعل نفسها في كل زمن. وللتعبير عن غياب زمن مخصوص، نكامل على جميع قيم \(t\):
\begin{equation}
\int_{-\infty}^{\infty}dt\,\psi^\dagger(0,t)\psi(0,t).
\end{equation}
وعند تأثير هذا المؤثر في الكيت الابتدائي، فإنه ينتج حالة. ويكون تداخلها مع البرا النهائي هو السعة للانتقال المحدد. ولا تُستخرج الاحتمالية إلا بعد أخذ مربع القيمة المطلقة للسعة.

ثمة عدد إضافي واحد نحتاج إليه. لِيَكُن \(g\) مقياسًا لشدة التفاعل بين الجسيم والهدف. ولأغراض الحساب في هذا النموذج، اكتب السعة القائدة التخطيطية على الصورة\footnote{تُحذَف عوامل تطبيع الحالات والعامل الكلي الاصطلاحي في نظرية الاضطراب من الرتبة الأولى؛ ولا يستخدم الحساب إلا قواعد اختيار المؤثرات والدعم الذي يفرض حفظ الطاقة.}
\begin{equation}
\mathcal A_{i\to f}
=g\left\langle k_f\left|
\int_{-\infty}^{\infty}dt\,
\psi^\dagger(0,t)\psi(0,t)
\right|k_i\right\rangle.
\end{equation}
يُسمّى المعامل \(g\) ثابت اقتران. وبالنسبة إلى هدف مشحون يشتت فوتونًا، توفّر الشحنة الكهربائية المقياس المناسب لشدة التفاعل: فزيادة الشحنة تزيد الميل إلى التشتت.

للتفاعلات المختلفة اقترانات مختلفة جدًا. فالميزون الذي يصل إلى بروتون يخضع لتفاعل قوي، وله شدة تشتت كبيرة نسبيًا. أما الفوتون الذي يشتته بروتون مشحون فله اقتران كهرومغناطيسي أصغر. وأما النيوترينو المتفاعل مع بروتون فله اقتران أصغر من ذلك أيضًا.

\begin{classroomqa}
\lecturetimestamp{01:09:09}
\audiencequestion{هل هذا النموذج البسيط صالح للإلكترون؟}
\lecturerresponse{فقط بوصفه بديلًا تبسيطيًا لعدة أنواع من التشتت. ولا ينبغي بعدُ أن يُفهم على أنه نظرية واقعية للإلكترون أو لأي هدف بعينه.}
\end{classroomqa}

سيكون التشتت الثقالي لغرافيتون عن نواة ضعيفًا على نحو استثنائي. ولا تشترك هذه الأمثلة في \(g\) عام واحد؛ ففي فيزياء الجسيمات ثوابت اقتران كثيرة.

ولا يلزم أيضًا أن يكون الاقتران الفعّال مستقلًا عن التردد على الدوام.

\begin{classroomqa}
\lecturetimestamp{01:10:31}
\audiencequestion{لماذا نأخذ \(g\) ثابتًا بدلًا من أن يكون دالة في \(\omega\) أو \(k\)؟}
\lecturerresponse{ليس من الضروري دائمًا أن يكون ثابتًا. فإذا امتصت ذرة فوتونًا، بقيت الذرة مثارة مدةً من الزمن، ثم يُصدَر الفوتون لاحقًا، وعندئذٍ يمكن أن يعتمد التشتت الفعّال على التردد. أما الثابت \(g\) فينتمي إلى النموذج الآني الحالي.}
\end{classroomqa}

في الوقت الحالي، أبقِ \(g\) ثابتًا واحسب نتائج هذا الاختيار الأبسط.

\begin{classroomqa}
\lecturetimestamp{01:12:32}
\audiencequestion{إذا كان \(g=1\)، فهل يجب أن يكون احتمال التشتت مساويًا لواحد؟}
\lecturerresponse{لا. فالتعبير الجاري حسابه ليس إلا التقريب القائد عند صِغَر \(g\). وإذا لم يكن الاقتران صغيرًا، لزم استعمال نظرية الاضطراب من رتب أعلى؛ وسيشير احتمال أكبر من واحد إلى انهيار تقريب الخطوة الواحدة، لا إلى احتمال فيزيائي.}
\end{classroomqa}

\section{اختزال مجاميع الأنماط}

عند \(x=0\)، تساوي الأطوار المكانية واحدًا، ومن ثم
\begin{align}
\psi(0,t)&=\sum_k a(k)e^{-i\omega(k)t},\\
\psi^\dagger(0,t)&=\sum_\ell a^\dagger(\ell)e^{i\omega(\ell)t}.
\end{align}

\begin{classroomqa}
\lecturetimestamp{01:13:58}
\audiencequestion{بما أن الهدف عند \(x=0\)، فهل ينبغي حذف الأسي المكاني؟}
\lecturerresponse{نعم. فعند الهدف \(e^{ikx}=e^{ik0}=1\)، كما أن العامل المكاني المرافق يساوي واحدًا أيضًا.}
\end{classroomqa}

استخدم رمزين صوريين مختلفين، \(k\) و\(\ell\)، في المجموعين. بالتعويض نحصل على
\begin{align}
\mathcal A_{i\to f}
&=g\int_{-\infty}^{\infty}dt
\sum_{\ell,k}
e^{i\omega(\ell)t}e^{-i\omega(k)t}
\langle k_f\rvert a^\dagger(\ell)a(k)\lvert k_i\rangle.
\end{align}
أي حد في مجموع \(k\) يمكن أن يبقى؟ يعطي مؤثر الإفناء صفراً في كل نمط خالٍ. تحتوي الحالة الابتدائية على جسيم في النمط \(k_i\) فقط، لذا \(k=k_i\). وأي حد في مجموع \(\ell\) يمكن أن يبقى؟ عند التأثير إلى اليسار، يجب أن يزيل \(a^\dagger(\ell)\) الجسيم الذي يمثله البرا النهائي، لذا \(\ell=k_f\). ينهار المجموعان معاً إلى
\begin{align}
\mathcal A_{i\to f}
&=g\int_{-\infty}^{\infty}dt\,
e^{i(\omega_f-\omega_i)t}
\langle k_f\rvert a^\dagger(k_f)a(k_i)\lvert k_i\rangle.
\end{align}
يحوّل مؤثر الإفناء كيت الجسيم الواحد الابتدائي إلى الفراغ،
\begin{equation}
a(k_i)\lvert k_i\rangle=\lvert0\rangle,
\end{equation}
ويحوّل مؤثر الخلق، عند تأثيره إلى اليسار، برا الجسيم الواحد النهائي إلى برا الفراغ،
\begin{equation}
\langle k_f\rvert a^\dagger(k_f)=\langle0\rvert.
\end{equation}
ولذلك يكون عنصر المصفوفة لهما
\begin{equation}
\langle0\vert0\rangle=1.
\end{equation}
ويُختزل كل تتبّع المؤثرات إلى واحد؛ ويعطي الاقتران \(g\)، والجزء غير البديهي المتبقي هو التكامل الزمني:
\begin{equation}
\mathcal A_{i\to f}
=g\int_{-\infty}^{\infty}dt\,
e^{i(\omega_f-\omega_i)t}.
\end{equation}

\editorialnote{سؤال قصير من الحضور عند 01:18:40 لا يمكن استعادته من المصدر. وتميّز الإجابة بين هذا الحساب ذي الجسيم الواحد وبين مسألة نسبية أو متعددة الجسيمات؛ ولا يمكن إعادة بناء تبادل أدق على نحو موثوق.}

الآن تنطبق متطابقة دالة دلتا السابقة مباشرة:
\begin{equation}
\int_{-\infty}^{\infty}dt\,
e^{i(\omega_f-\omega_i)t}
=2\pi\delta(\omega_f-\omega_i).
\end{equation}
ومن ثم
\begin{equation}
\boxed{
\mathcal A_{i\to f}
=2\pi g\,\delta(\omega_f-\omega_i)
}.
\end{equation}
تنعدم السعة ما لم
\begin{equation}
\omega_f=\omega_i.
\end{equation}
وبما أن \(\hbar=1\)، فهذا هو \(E_f=E_i\). وقد ظهر حفظ الطاقة لأن كل أزمنة التفاعل عوملت على السواء. ولو قُيِّد التكامل الزمني إلى فترة منتهية لأنتج دالة متسعة بدلاً من دالة دلتا مضبوطة تماماً. إن اللاتغير الدقيق تحت انتقالات الزمن هو ما يعطي شرط حفظ الطاقة الدقيق.

\begin{classroomqa}
\lecturetimestamp{01:21:30}
\audiencequestion{يجب تربيع السعة للحصول على احتمال. ماذا يحدث لدالة دلتا عندما تُربَّع؟}
\lecturerresponse{تتطلب دالة دلتا المربعة عناية إضافية. والنقطة المطلوبة هنا مؤكدة: فالنتيجة تساوي صفراً ما لم تتساوَ الطاقتان الابتدائية والنهائية. وعلى الحامل الذي يتحقق عليه حفظ الطاقة، يحتوي العامل المضاعِف التربيعي التخطيطي على \(4\pi^2g^2\).}
\end{classroomqa}

\editorialnote{لا يمكن الحصول على احتمال انتقال أو معدل انتقال مطبَّع بمجرد طرح \(\delta(\omega_f-\omega_i)^2\). فذلك يتطلب تنظيماً بزمن منتهٍ وتطبيعاً متسقاً للحالات. لذلك لا يُحتفَظ بـ \(4\pi^2g^2\) إلا بوصفه العامل المضاعِف التربيعي التخطيطي الذي شدد عليه المصدر.}

لا يوجد في تفاعل الهدف النقطي هذا ما يميز زخماً نهائياً عن آخر سوى شرط الطاقة. لذلك، في ثلاثة أبعاد، يعطي النموذج، عند الرتبة الرائدة، الوزن نفسه لجميع الزخوم الخارجة ذات الطاقة نفسها---وبصورة مكافئة، لجميع اتجاهات الخروج على الكرة ذات الطاقة الثابتة. إن هذا التساوي الخواصي خاصية لهذا المبعثر النقطي البسيط جداً، لا لكل هدف حقيقي.

قدّم الحساب ثلاث أفكار مترابطة. فالمعامل \(g\) يقيس شدة التفاعل. والتعبير من الشكل برا-مؤثر-كيت هو سعة تشتت يدخل مربع قيمتها المطلقة في الاحتمال. والأهم من ذلك أن التكامل على جميع الأزمنة ينتج دالة دلتا للطاقة: فتناظر انتقال الزمن يفرض حفظ الطاقة.

\section{ما الذي يتغير عندما يتغير الهدف أو الاقتران؟}

اختير الأصل \(x=0\) لمجرد التيسير. ضع الهدف عند \(x=x_{\mathrm T}\). تعطي العوامل المكانية في \(\psi^\dagger(x_{\mathrm T},t)\psi(x_{\mathrm T},t)\)
\begin{equation}
e^{-ik_f x_{\mathrm T}}e^{ik_i x_{\mathrm T}}
=e^{i(k_i-k_f)x_{\mathrm T}}.
\end{equation}
ومن ثم تصبح السعة
\begin{equation}
\mathcal A_{i\to f}(x_{\mathrm T})
=e^{i(k_i-k_f)x_{\mathrm T}}
\mathcal A_{i\to f}(0).
\end{equation}

\begin{classroomqa}
\lecturetimestamp{01:24:37}
\audiencequestion{هل يمكن دائماً إعادة تعريف الموضع بحيث يكون الهدف عند \(x=0\)؟}
\lecturerresponse{بالنسبة إلى هذا الهدف الثابت الوحيد، نعم. وإذا تُرك عند \(x_{\mathrm T}\neq0\)، فإن موضعه يظهر على هيئة الطور \(e^{i(k_i-k_f)x_{\mathrm T}}\) في السعة. والضرب في الطور المرافق عقدياً يزيله من الاحتمال.}
\end{classroomqa}

في الواقع،
\begin{equation}
\left|e^{i(k_i-k_f)x_{\mathrm T}}\right|^2=1.
\end{equation}
تحتفظ السعة بأثر موضع الهدف، لكن احتمال الهدف الواحد هذا لا يحتفظ به.

والآن اعكس هذا اللاتناظر. افترض أن التفاعل يمكن أن يحدث عند أي موضع، ولكن عند زمن مميز واحد فقط. عندئذٍ يعطي التكامل على الموضع
\begin{equation}
\int_{-\infty}^{\infty}dx\,
e^{i(k_i-k_f)x}
=2\pi\delta(k_i-k_f),
\end{equation}
في حين أن تثبيت الزمن يزيل دالة دلتا للطاقة.

\begin{classroomqa}
\lecturetimestamp{01:26:41}
\audiencequestion{إذا لم يكن مهماً أين حدث التفاعل، ولكنه حدث عند زمن ثابت واحد فقط، فهل تُحفَظ كمية مختلفة؟}
\lecturerresponse{نعم. فالتكامل على جميع المواضع يفرض حفظ الزخم، بينما الزمن المميز للتفاعل يزيل حفظ الطاقة الدقيق. وفي مسألة كاملة لجسيمين، قد يقع التفاعل في أي زمان وأي مكان.}
\end{classroomqa}

وهذا يعطي أيضاً تفسيراً فيزيائياً للاعتماد الزمني الصريح. فإذا كان الاقتران هو \(g(t)\)، أو إذا جعل تدخل خارجي التشتت ممكناً فجأة، فليست كل الأزمنة متكافئة. وعندئذ قد يتبادل الوصف المختزل الطاقة مع أي ديناميات خارجية أُغفلت.

\begin{classroomqa}
\lecturetimestamp{01:28:29}
\audiencequestion{هل تأتي ثوابت الاقتران في الغالب من التجربة، وما معنى اقتران يعتمد على الزمن؟}
\lecturerresponse{نعم، فقيمها تُحدَّد تجريبياً. والاعتماد الزمني الصريح يعني عادةً أن جزءاً متحركاً أو متغيراً من النظام قد تُرك خارج الوصف، تماماً كما تُرك ارتداد الهدف خارج نموذج الهدف الثابت. والطاقة التي تبدو مفقودة تعود إلى درجات الحرية المُغفلة تلك.}
\end{classroomqa}

الهدف المتحرك هو أبسط مثال.

\begin{classroomqa}
\lecturetimestamp{01:29:23}
\audiencequestion{هل إصابة هدف متحرك مثال لا تُحفَظ فيه طاقة الجسيم؟}
\lecturerresponse{نعم. فإذا أُهملت الحركة المفروضة للهدف من حساب الطاقة، أمكنها أن تزود الجسيم بالطاقة أو أن تسحبها منه.}
\end{classroomqa}

الطابع النقطي للهدف تبسيط آخر.

\begin{classroomqa}
\lecturetimestamp{01:29:43}
\audiencequestion{متى يجب مقارنة حجم الهدف بحجم الموجة الساقطة أو بطولها الموجي؟}
\lecturerresponse{في النموذج الحالي، حجم الهدف صفري. والحجم المنتهي للهدف يولد بالفعل آثاراً إضافية عند مقارنته بالطول الموجي، لكن هذا التنقيح مؤجل.}
\end{classroomqa}

عُد الآن إلى الهدف المتحرك. تخيل جداراً يتحرك نحو كرة تنس. يمكن للجدار أن يعطي الكرة طاقة. وإذا كانت حركة الجدار مفروضة وأُهملت طاقته، فإن طاقة الكرة وحدها لا تُحفَظ. رياضياً، يكون هذا اعتماداً زمنياً في موضع الهدف لا بالضرورة في \(g\).

ومن ثم فإن مؤثرات الخلق والإفناء قد أدت بالضبط ما أُريد لها: فهي تصف انتقالاً من حالة جسيم واحد إلى أخرى، بينما يحدد الاعتماد الزمكاني للحقول شروط الانحفاظ التي ترافق ذلك الانتقال.

\section{عدّ الشحنة وثبات الطور}

لكل نوع من الجسيمات حقله الخاص. فالإلكترونات والفوتونات، على سبيل المثال، لا يصفها \(\psi\) نفسه. ولإبراز عدّ الشحنة بالجبر المتاح بالفعل، لندع \(\psi\) يرمز مؤقتاً إلى حقل شبيه بالإلكترون: يفني \(\psi\) جسيمًا مشحونًا واحدًا من هذا النوع، ويخلق \(\psi^\dagger\) واحدًا. وهذا بديل بوزوني لغرض ضبط الحساب؛ أما حقل الإلكترون الفيرميوني الحقيقي فيتطلب جبراً لم يُدخَل بعد.

العامل الخاص بالهدف الثابت
\begin{equation}
\psi^\dagger\psi
\end{equation}
يفني جسيمًا واحدًا ويخلق جسيمًا واحدًا. يدخل إلكترون واحد ويخرج إلكترون واحد، لذا لا تتغير الشحنة التي يحملها هذا النوع.

تخيل الآن دخول إلكترون واحد وخروج إلكترونين من النقطة نفسها. سيحتوي جداء المؤثرات الموافق على حقل إفناء واحد وحقلين للخلق،
\begin{equation}
\psi^\dagger\psi^\dagger\psi.
\end{equation}
وبالنسبة إلى حقل إلكترون، فإن هذا سينشئ إلكترونًا صافياً واحدًا وينتهك حفظ الشحنة الكهربائية. ويمكن بدلاً من ذلك تمثيل دخول إلكترونين وخروج إلكترونين بـ
\begin{equation}
\psi^\dagger\psi^\dagger\psi\psi,
\end{equation}
إذ يحتوي هذا على أعداد متساوية من حقول الخلق والإفناء. أما دخول اثنين وخروج ثلاثة فسيتطلب
\begin{equation}
\psi^\dagger\psi^\dagger\psi^\dagger\psi\psi,
\end{equation}
ويغير مرة أخرى الشحنة الصافية. والقاعدة المباشرة بسيطة: فالتفاعل الحافظ للشحنة يحتوي العدد نفسه من \(\psi\) و\(\psi^\dagger\).

\begin{classroomqa}
\lecturetimestamp{01:34:27}
\audiencequestion{هل يهم ترتيب مؤثرات الحقل في اختبار حفظ الشحنة هذا؟}
\lecturerresponse{الاختبار المباشر يتعلق بعدد عوامل \(\psi\) و\(\psi^\dagger\)، لا بترتيبها. وأما ما إذا كانت الجداءات المختلفة الترتيب متطابقة في غير ذلك، فذلك سؤال منفصل.}
\end{classroomqa}

هناك معيار رياضي أعمق من مجرد العد. اضرب الحقل في طور ثابت:
\begin{equation}
\psi\longrightarrow e^{i\alpha}\psi.
\end{equation}
وبما أن \(\psi^\dagger\) هو المرافق الهيرميتي، فإنه يتحول بالطور المرافق،
\begin{equation}
\psi^\dagger\longrightarrow e^{-i\alpha}\psi^\dagger.
\end{equation}
والتعبير الذي يحتوي أعداداً متساوية من الحقلين يبقى دون تغير. فعلى سبيل المثال،
\begin{align}
\psi^\dagger\psi
&\longrightarrow
e^{-i\alpha}e^{i\alpha}\psi^\dagger\psi
=\psi^\dagger\psi,\\
\psi^\dagger\psi^\dagger\psi\psi
&\longrightarrow
e^{-2i\alpha}e^{2i\alpha}
\psi^\dagger\psi^\dagger\psi\psi
=\psi^\dagger\psi^\dagger\psi\psi.
\end{align}
أما التعبير الذي يحتوي أعداداً غير متساوية فيتغير. ومن ثم
\begin{equation}
\psi^\dagger\psi^\dagger\psi
\longrightarrow
e^{-i\alpha}\psi^\dagger\psi^\dagger\psi.
\end{equation}
فإن التفاعلات المسموح بها الحافظة للشحنة هي بالضبط تلك الثابتة تحت هذا التغير العالمي في الطور. وثبات الطور وحفظ الشحنة وصفان للقاعدة نفسها في لغة الحقول البسيطة هذه.

التركيب
\begin{equation}
\psi+\psi^\dagger
\end{equation}
لا يبقى ثابتاً تحت دوران طور اعتباطي.

\begin{classroomqa}
\lecturetimestamp{01:37:53}
\audiencequestion{أليس ذلك هو الجزء الحقيقي من \(\psi\)؟}
\lecturerresponse{بلى. إنه يتناسب مع الجزء الحقيقي من \(\psi\)، وضرب \(\psi\) بعامل طور يغيّر ذلك الجزء الحقيقي. لذلك فهو ليس ثابتًا تحت تحولات الطور.}
\end{classroomqa}

\begin{classroomqa}
\lecturetimestamp{01:38:05}
\audiencequestion{إذا أصاب إلكترونٌ هدفًا وانبعثت إلكترونات ثانوية كثيرة، فكيف يمكن أن تبقى الشحنة محفوظة؟}
\lecturerresponse{ليس الهدف خاملًا في تلك العملية. فإذا أصدرت ذرة إلكترونات كثيرة، فإن شحنة الذرة نفسها وحالتها تتغيران. ويجب إدراج ديناميكيتها، كما يجب أيضًا وصف الهدف بمجالات كمومية.}
\end{classroomqa}

عند هذه النقطة تظهر ثلاث مراسلات. تماثل الانسحاب الزمني يقابل حفظ الطاقة. وثبات الطور العام يقابل حفظ الشحنة. وتماثل الانسحاب المكاني يقابل حفظ الزخم، مع أن الحساب المكاني الكامل لا يزال يتعيّن إنجازه.

\section{ترتيب المؤثرات وحدود نموذج أدنى رتبة}

تبسيطًا، تضع الجداءات أعلاه جميع حقول الإنشاء إلى اليسار وجميع حقول الإفناء إلى اليمين. ويمكن لجداءٍ متداخل أن يظل ذا أعداد متساوية، ومن ثم يجتاز اختبار عدّ الشحنة، لكن ترتيب المؤثرات قد يؤثر في التعبير بطرائق أخرى.

\begin{classroomqa}
\lecturetimestamp{01:40:02}
\audiencequestion{هل يمكن كتابة النوع نفسه من العملية بعوامل متداخلة مثل \(\psi\psi^\dagger\psi\psi^\dagger\)؟}
\lecturerresponse{ما زال يحوي أعدادًا متساوية، ولذلك يحقق قاعدة الشحنة، لكنه ليس عمومًا التعبير نفسه كتعبيرٍ تكون فيه جميع الخناجر إلى اليسار. ويؤجَّل تناول إعادة ترتيب هذه التراتيب المختلفة.}
\end{classroomqa}

\editorialnote{صياغة سؤالٍ من الجمهور عند 01:41:10 متضررة أكثر مما يسمح بإعادة بنائها. ويكرر الرد الذي أمكن استخلاصه أن تساوي عددَي \(\psi\) و\(\psi^\dagger\) يستلزم ثبات الطور وحفظ الشحنة، في حين أن عدم تساويهما يعني غياب الأمرين.}

هذه الروابط بين التناظرات وقوانين الانحفاظ هي أمثلة على مبرهنة نويثر. ففي الميكانيكا الكلاسيكية تمنح تلك المبرهنة الإطار العام؛ وهنا عُرض الارتباط نفسه مباشرةً من خلال الأطوار ودوال دلتا الخاصة بالحقول الكمومية.

للنتيجة عند أدنى رتبة في حالة الهدف الثابت أيضًا مجال صلاحية واضح. ولا يمكن الوثوق بمعامل أمامي احتمالي تخطيطي مثل \(4\pi^2g^2\) إذا اقترب من الواحد أو تجاوزه.

\begin{classroomqa}
\lecturetimestamp{01:43:33}
\audiencequestion{هل تعني النتيجة التخطيطية أن \(g\) يجب أن يكون أصغر تقريبًا من \(1/(2\pi)\)؟}
\lecturerresponse{إما أن يكون \(g\) صغيرًا بما يكفي لكي يكون هذا الحد الأول مفيدًا، وإما أن يكون النموذج غير مكتمل. يمكن لجسيم أن يتشتت مرة، ثم يعود فيتشتت مرتين، ثم ثلاث مرات، وهكذا. والسعة الكاملة سلسلة لا نهائية. ولا يهيمن أبسط حد إلا عندما يكون \(g\) صغيرًا؛ وإلا وجب جمع كثير من مخططات فاينمان.}
\end{classroomqa}

هذه الصور البسيطة للزمكان هي بدايات مخططات فاينمان. وتحدد رؤوسها الإمكانات الأولية: دخول جسيم واحد وخروج جسيم واحد يقابل عاملًا مثل \(\psi^\dagger\psi\)؛ ودخول جسيمين وخروج جسيمين يقابل عاملًا مثل \(\psi^\dagger\psi^\dagger\psi\psi\). وتولد الرؤوس المتكررة قوى أعلى لثابت الاقتران وحدودًا من رتب أعلى في السلسلة.

\begin{classroomqa}
\lecturetimestamp{01:45:39}
\audiencequestion{هل يمكن أن يكون لثابت اقتران مركبة تخيلية؟}
\lecturerresponse{يمكن أن تظهر ثوابت اقتران مركبة. ويمكن لتركيبات الطور الفردية تحت \(T\) وغير القابلة للاختزال أن تشير إلى انتهاك تناظر عكس الزمن: فالعملية المعكوسة زمنيًا قد لا تحمل السعة المرتبطة بالتناظر.}
\end{classroomqa}

\editorialnote{ذُكر معامل \(\theta\) في الديناميكا اللونية الكمومية في هذا التبادل، لكنه معامل حقيقي يضرب تآثرًا فرديًا تحت \(T\)، لا مركبةً تخيلية لثابت اقتران. والطور القابل للإزالة أو الطور المركب الامتصاصي لا يثبت بمفرده انتهاك تناظر \(T\). ولا ينطبق هذا الاستنتاج إلا على تركيبات الطور الفردية تحت \(T\) وغير القابلة للاختزال.}

\begin{classroomqa}
\lecturetimestamp{01:47:02}
\audiencequestion{هل يمكن أن ينبعث الجسيم الخارج قبل امتصاص الجسيم الداخل أو بعده، بدلًا من حدوث الأمرين في اللحظة نفسها؟}
\lecturerresponse{هذا سؤال وجيه، لكن النموذج الآني لا يتناوله. ويؤجَّل الترتيب الزمني الأعم.}
\end{classroomqa}

\begin{classroomqa}
\lecturetimestamp{01:47:44}
\audiencequestion{كيف يُصاغ هذا البناء بصيغة نسبية؟}
\lecturerresponse{الإشارة الموجزة هي جمع الأجزاء المرافقة وإدراج مقام يتضمن جذرًا تربيعيًا لـ \(\omega(k)\). ويؤجَّل البناء الكامل.}
\end{classroomqa}

\editorialnote{الصيغة النسبية الموجزة التي وردت شفهيًا لا يمكن استعادتها على نحو موثوق بما يكفي لتنضيدها. ولم يُحتفظ إلا بمكوّناتها المذكورة وبالتأجيل المشار إليه.}

ويجب أيضًا التمييز بوضوح بين حفظ الشحنة وحفظ الكتلة الساكنة.

\begin{classroomqa}
\lecturetimestamp{01:48:22}
\audiencequestion{هل تُحفَظ الكتلة الساكنة بالمعنى البنيوي نفسه الذي تُحفَظ به الشحنة الكهربائية؟}
\lecturerresponse{لا. في الميكانيكا الجاليليّة غير النسبية، تكون الكتلة كمية جمعية محفوظة. أما في التفاعلات النسبية، فليس من الضروري أن يكون مجموع الكتل الساكنة للجسيمات المفردة محفوظًا. ويمكن لإلكترون وبوزيترون أن يفنيا إلى فوتونين مع بقاء الشحنة الكهربائية محفوظة.}
\end{classroomqa}

\editorialnote{في نظام نسبي كامل ومغلق، يُحفَظ الزخم الرباعي الكلي، فتكون كتلته اللامتغيرة ثابتة. ويتعلق مثال الإفناء بمجموع الكتل الساكنة للجسيمات المفردة، لا بالكتلة اللامتغيرة للنظام المغلق.}

وعليه فإن حفظ الكتلة الجمعية حقيقة فعالة غير نسبية، لا قاعدة أساسية على قدم المساواة مع حفظ الشحنة الكهربائية. وثبات الطور الذي وُجد أعلاه يعبّر عن تناظر الشحنة الأعمق.

والخطوات التالية هي استبدال البديل البوزوني بحقول فرميونية، والسماح بعدة أنواع من الجسيمات---الإلكترونات والميونات والكواركات والفوتونات وغيرها---لكلٍّ منها حقله الخاص وتآثراته المسموح بها.

\begin{classroomqa}
\lecturetimestamp{01:50:39}
\audiencequestion{هل كان من المشروع أن يُسمّى المثال السابق حقلًا إلكترونيًا مع أن الحقول البوزونية وحدها هي التي طُوِّرت حتى الآن؟}
\lecturerresponse{فقط بوصفه نسخة بوزونية أو بديلًا عن الإلكترون. أما النظرية الفرميونية الحقيقية فلم تُقدَّم بعد.}
\end{classroomqa}